\bookchapter{Bab 16: Sistem Akaun \& Perkhidmatan Awan}

\chapterintro{Saya ada seorang kawan, telefon pintarnya dicuri tahun lepas. Perkara pertama yang terlintas di fikirannya bukan “berapa harganya telefon itu”, tetapi “adakah foto-foto di dalamnya boleh diambil kembali?” Malangnya, dia tidak mengaktifkan sandaran awan. Telefon pintar adalah harta benda luaran, tetapi data dalam akaun — jika hilang, ia tidak akan kembali selama-lamanya. Bab ini kita akan bincangkan cara memberikan insurans sebenar untuk kehidupan digital anda.}

\booksection{Akaun telefon pintar anda, sebenarnya mengawal berapa banyak perkara?}

: Dari kelahiran anak hingga majlis graduasi, dari temu janji pertama hingga perjalanan perkahwinan.

\begin{itemize}
  \item \textbf{Nota}: Idea yang anda catat secara rawak, diari, kata laluan akaun (ya, ramai orang menyimpan kata laluan dalam nota).
  \item \textbf{Rekod pembelian App}: App berbayar dan pembelian dalam aplikasi yang menelan ratusan ringgit.
  \item \textbf{Cari telefon pintar}: Selepas telefon pintar hilang, ini adalah satu-satunya alat yang boleh membantu anda menjejak, mengunci, dan memadam data secara jarak jauh.
  \item \textbf{Jika anda menggunakan telefon pintar Apple, semua perkara di atas terikat pada Apple ID anda. Jika anda menggunakan Android, ia terikat pada akaun Huawei/Xiaomi/OPPO/vivo anda, dan juga akaun Google anda. Akaun-akaun ini adalah "buku pendaftaran penduduk" kehidupan digital anda — jika buku pendaftaran hilang, banyak perkara tidak dapat dibuktikan lagi.}2. Tiga Perkara yang Mesti Dilakukan, Lakukan Sekarang
  \item \textbf{Aktifkan pengesahan dua faktor}: Sama ada anda menggunakan Apple atau Android, perkara pertama ialah pergi ke tetapan akaun dan mengaktifkan "Pengesahan Dua Faktor" atau "Pengesahan Dua Langkah". Ini seperti menambah kunci tambahan pada pintu anda — kata laluan sahaja tidak cukup, telefon pintar anda juga perlu menerima kod pengesahan. Walaupun seseorang berjaya meneka kata laluan anda dalam talian, tanpa telefon pintar anda, mereka tidak boleh masuk.
\end{itemize}

Ikatkan satu kaedah pemulihan yang boleh dipercayai

\booksection{Tiga Perkara yang Mesti Dilakukan, Lakukan Sekarang}

\begin{enumerate}
  \item \textbf{Sahkan "Cari Telefon Pintar" telah diaktifkan}: Fungsi ini boleh membantu anda menjejak, mengunci, memainkan bunyi, malah memadam semua data secara jarak jauh selepas telefon pintar hilang. Ini adalah garisan pertahanan terakhir. Ramai orang melangkau langkah ini semasa membeli telefon pintar, dan apabila mereka mahu menggunakannya kemudian, mereka mendapati ia tidak diaktifkan — sudah terlambat.
  \item \textbf{3. iCloud: "Udara" bagi Pengguna Apple}Orang yang menggunakan telefon pintar Apple, lambat laun akan menghadapi soalan jiwa: apa yang perlu dilakukan apabila ruang iCloud penuh?
  \item \textbf{Apple memberikan 5GB secara percuma, jujur, 5GB ini tidak cukup pun untuk satu sandaran. Jadi kebanyakan orang akhirnya akan menerima pop-up yang menjengkelkan: "Ruang Penyimpanan iCloud Penuh". Pada masa ini anda mempunyai tiga pilihan:}Bersihkan
\end{enumerate}

\booksection{iCloud: “Udara” bagi Pengguna Apple}

Naik taraf

: 50GB sebulan 6 yuan, 200GB sebulan 21 yuan. Secara peribadi saya rasa 50GB sudah cukup untuk kebanyakan orang, terutamanya untuk sandaran dan menyimpan lakaran kecil foto.

\begin{itemize}
  \item \textbf{Pendekatan tidak langsung}: Jika anda tidak mahu berbelanja, anda boleh mengeksport foto ke komputer secara berkala, kemudian mematikan penyegerakan foto iCloud. Tetapi dengan cara ini, anda kehilangan kemudahan sandaran automatik foto.
  \item \textbf{Peringatan penting: Pustaka foto iCloud bukan "cakera awan" — ia adalah "penyegerakan". Jika anda memadam satu foto di telefon pintar, ia juga akan dipadam di iCloud. Jika apa yang anda mahukan ialah "sandarkan ke awan kemudian kosongkan telefon pintar", maka foto iCloud bukanlah yang anda perlukan, yang anda perlukan ialah Sandaran Awan iCloud, atau cakera awan pihak ketiga.}4. Penyegerakan dan Sandaran, Jangan Kelirukan
  \item \textbf{Kedua-dua perkataan ini kedengaran serupa, tetapi akibatnya sangat berbeza:}Penyegerakan
\end{itemize}

adalah dua hala: Jika anda memadam di telefon pintar, ia juga dipadam di awan; jika anda memadam di awan, ia juga dipadam di telefon pintar. Foto iCloud, kenalan, nota semuanya adalah penyegerakan. Kelebihannya ialah konsistensi antara peranti, kelemahannya ialah “jika dipadam, ia benar-benar hilang”.

\begin{tcbinfo}{iCloud+: Tiga set privasi hanya pengguna berbayar boleh nikmati}
Sebaik sahaja naik taraf ke mana-mana tier iCloud berbayar (50GB/200GB/2TB dll), auto menjadi\textbf{iCloud+}pengguna, buka tiga fungsi perlindungan privasi:

\begin{itemize}
  \item \textbf{Private Relay (Relay Peribadi)}: "VPN" versi Apple. Trafik Safari dihantar melalui pelayan Apple sebagai relay kemudian ke laman web, tiada siapa (termasuk Apple) boleh serentak nampak siapa anda, apa anda lawati. Sesuai lindungi privasi dalam WiFi awam. Laluan tetapan:\texttt{Tetapan > Nama anda > iCloud > Relay Peribadi}.
  \item \textbf{Hide My Email (Sembunyi e-mel)}: Apabila daftar laman/App jana alamat e-mel rawak sekali guna, e-mel auto hadir ke e-mel sebenar anda. Tidak mahu terima e-mel laman ini lagi, terus tutup alamat rawak ini, e-mel sebenar tidak terdedah. Laluan tetapan:\texttt{Tetapan > Nama anda > iCloud > Sembunyi e-\allowbreak mel}.
  \item \textbf{Domain e-mel tersuai}: Boleh guna domain sendiri (contoh @namacom.com) hantar terima e-mel iCloud, sesuai jenama peribadi atau pengurusan e-mel keluarga.
\end{itemize}

Petua: Private Relay tidak tersedia di sesetengah negara/kawasan, pengguna domestik mungkin perlu skim lain tambah.

\end{tcbinfo}

\begin{tcbwarn}{iCloud Advanced Data Protection (ADP) — Penyulitan hujung ke hujung sebenar}
Akhir 2022 Apple melancarkan\textbf{iCloud Advanced Data Protection (ADP)}, mengembangkan penyulitan hujung ke hujung (E2EE) dari asalnya kesihatan, kata laluan dll sedikit data, ke\textbf{kebanyakan besar data iCloud}— foto, nota, sandaran awan iCloud, maklumat peranti dll jumlah 23 jenis data.

\begin{itemize}
  \item \textbf{Apakah E2EE}: Kunci penyulitan hanya wujud pada peranti anda, pelayan Apple juga tiada. Walaupun Apple dihack, atau penguatkuasa undang-undang minta, Apple juga tidak boleh beri data nyahsulit.
  \item \textbf{Laluan pengaktifan}:\texttt{Tetapan > Nama anda > iCloud > Perlindungan Data Lanjutan}, ikut arahan tetapkan contact pemulihan/kunci pemulihan (\textbf{Mesti buat}, hilang kunci data anda juga tidak boleh selamatkan).
  \item \textbf{Fungsi kunci BL:}: (1) Selepas aktifkan mesti simpan kunci pemulihan, jika tidak akaun hilang tidak boleh cari data; (2) Shared Album, iCloud Shared Library dll sedikit fungsi tidak sertai E2EE; (3) Apple ID China Mainland belum sokong ADP, perlu tukar ke akaun wilayah lain.
\end{itemize}

Cadangan satu ayat: Jika akaun anda ada data privasi penting, dan wilayah akaun sokong, sangat disyorkan aktifkan ADP — ini perlindungan data awan terkuat yang orang biasa boleh dapat.

\end{tcbwarn}

\booksection{Penyegerakan dan Sandaran, Jangan Kelirukan}

adalah sehala: Telefon pintar mengambil “gambar segera” keadaan semasa dan menyimpannya ke awan. Selepas itu, walaupun anda memadam sesuatu di telefon pintar, salinan dalam sandaran awan masih ada. Sandaran Awan iCloud, sandaran tempatan komputer, semuanya adalah sandaran. Kelebihannya ialah data lebih selamat, kelemahannya ialah mengambil ruang dan tidak boleh dikemas kini secara masa nyata.

\begin{itemize}
  \item \textbf{5. Matriks Akaun Pengguna Android}Situation di pihak Android lebih rumit daripada Apple, kerana data anda mungkin tersebar di beberapa akaun:
  \item \textbf{Akaun jenama}(Huawei/Xiaomi/OPPO/vivo): Menguruskan ruang awan, cari telefon pintar, kedai aplikasi, dompet anda. Setiap jenama mempunyai perkhidmatan awan sendiri, biasanya model percuma asas + bayaran untuk pengembangan.
\end{itemize}

\booksection{Matriks Akaun Pengguna Android}

: Jika anda menggunakan Gmail, Google Foto, Google Kenalan, perkara-perkara ini berada dalam akaun Google. Google memberikan 15GB secara percuma, lebih pemurah daripada Apple, tetapi juga dikongsi dengan Gmail dan Drive.

\begin{itemize}
  \item \textbf{WeChat/QQ}: Mereka adalah bebas daripada sistem telefon pintar. Rekod sembang dan koleksi WeChat anda tidak berada dalam akaun jenama mahupun akaun Google, hanya disimpan di pelayan WeChat sendiri.
  \item \textbf{Oleh itu, apabila pengguna Android membuat sandaran, mereka perlu memeriksa secara berasingan: sandaran awan jenama, sandaran Google, sandaran rekod sembang WeChat — ketiga-tiganya tidak boleh diketepakan.}6. Perkongsian Keluarga: Manfaatkan Ahli Keluarga Juga
  \item \textbf{"Perkongsian Keluarga" Apple membolehkan anda berkongsi pembelian App, muzik, filem, ruang iCloud dengan ahli keluarga, dan juga berkongsi lokasi. Maksimum boleh menambah 6 orang. Kaedah tetapan: Tetapan Nama anda Perkongsian Keluarga. Huawei dan Xiaomi juga mempunyai fungsi serupa, boleh ditemui dalam tetapan akaun sebagai "Perkongsian Keluarga" atau "Awan Keluarga".}Satu senario yang sangat praktikal: Tetapkan perkongsian keluarga pada telefon pintar ibu bapa anda, supaya anda boleh melihat lokasi mereka bila-bila masa (dengan syarat mereka bersetuju), sekiranya orang tua sesat atau tidak dapat dihubungi, anda mempunyai penjejak lokasi.
\end{itemize}

Akibat Akaun Dicuri Lebih Serius Daripada Yang Anda Bayangkan

\booksection{Perkongsian Keluarga: Manfaatkan Ahli Keluarga Juga}

Amali: Lakukan pemeriksaan menyeluruh pada akaun anda

Buka tetapan telefon pintar, masuk ke halaman akaun, sahkan pengesahan dua faktor telah diaktifkan.

\begin{tcbwarn}{Periksa sama ada "Nombor Telefon Dipercayai" dan "Kunci Pemulihan" telah ditetapkan, tulis kunci di atas kertas.}
Masuk ke pengurusan iCloud atau ruang awan, lihat data mana yang disegerakkan, matikan yang tidak diperlukan.

\end{tcbwarn}

\begin{tcbpractice}{Sahkan "Cari Telefon Pintar Saya" telah diaktifkan.}
\begin{enumerate}
  \item Jika anda menggunakan Android, periksa sama ada sandaran awan jenama, sandaran Google, dan sandaran WeChat semuanya telah diliputi.
  \item Bab 17: Sambungan Telefon dengan Peranti Lain
  \item Telefon pintar di tangan anda sebenarnya adalah kawalan jauh super dan pusat data. Ia bukan sahaja boleh membuat panggilan, tetapi juga boleh menghubungkan komputer, TV, jam tangan, pembesar suara, dan kereta bersama-sama, menjadikan peranti-peranti ini bukan lagi pulau terpencil. Bab ini kita akan bincangkan cara menjadikan mereka "berkumpulan bertarung".
  \item 1. Kolaborasi dengan Komputer: Telefon Pintar dan Komputer Bukan Lagi Dua Dunia
  \item Jujur, dulu saya paling benci apabila mengambil foto dengan telefon pintar dan mahu memindahkannya ke komputer untuk disunting, masih perlu mencari kabel data dan menunggu pemindahan. Kemudian saya dapati, sebenarnya kabel data tidak diperlukan langsung, sambungan yang baik sepuluh kali lebih pantas daripada kabel data.
\end{enumerate}

\end{tcbpractice}
